% docs/.paper/sec_conclusion.tex
\section{Conclusion}
\label{sec:conclusion}

This paper introduced \textbf{MAgHARCM}, a modular multi-agent harness for repository-scale code translation engineered to run entirely on local Small Language Models (SLMs) served by Ollama. MAgHARCM coordinates four specialised agents---Analyzer, Planning, Translator, and Validator---over an explicit \textsc{Analyze} $\to$ \textsc{Plan} $\to$ \textsc{Translate} $\to$ \textsc{Validate} $\to$ \textsc{Repair} loop. Cycle-free reverse-topological dependency scheduling with back-edge removal enforces callee-before-caller ordering; a $\le 4$\,KB Symbol Navigator and a $\le 4$\,KB Prior-Modules Memory bound per-fragment context; an adversarial test-weakening guard prevents repair loops from inflating pass rates by removing assertions; and a CodaMOSA-style coverage-plateau detector bounds repair spinning.

Our empirical evaluation across four open-source benchmark repositories spanning C, Go, and Java targeting Rust shows that local pairings of a 30B reasoning model with an instruction-tuned 4B coding model achieve clean compilation on every project without cloud dependencies or data exposure. The harness converges to a 100\% test pass rate on algorithmic code (GildedRose), 62.5\% on streaming libraries (Gohistogram), 57.1\% on statistical suites (Stats), and 26.5\% on enterprise validation frameworks (Apache Commons Validator) before the plateau detector terminates the repair loop. Our failure taxonomy identifies borrow-checker lifetime errors (38\%) and cyclic package references (26\%) as the dominant bottlenecks for local-SLM repository migration.

MAgHARCM is open-source and reproducible; the artifacts (harness, benchmark configurations, agent logs, and per-iteration checkpoints) accompany this submission. Future work will extend the target language beyond Rust, integrate the optional ABCoder MCP provider for richer cross-language structural queries, and evaluate the four-agent decomposition on multi-PL-pairs (e.g.\ Go-to-Rust \emph{and} Java-to-Python) using a single harness configuration.
